\documentclass[11pt]{article}

\usepackage[margin=1in]{geometry}
\usepackage{fontspec}
\usepackage{microtype}
\usepackage{enumitem}
\usepackage{xcolor}
\usepackage{titlesec}
\usepackage[colorlinks=true,linkcolor=black,urlcolor=blue]{hyperref}

% Monospace for inline code; \code{} keeps underscores/braces literal.
\newcommand{\code}[1]{\texttt{#1}}

% Blockquote environment for the two pull-quotes.
\newenvironment{pullquote}
  {\begin{quote}\itshape}
  {\end{quote}}

\titleformat{\section}{\Large\bfseries}{}{0em}{}
\titleformat{\subsection}{\large\bfseries}{}{0em}{}

\setlist[itemize]{leftmargin=1.5em,itemsep=0.3em}

\title{Case study: security is a second-class citizen in pi}
\author{Stacia}
\date{}

\begin{document}
\maketitle

\noindent An attempt to understand pi's permission model surfaced a structural
flaw in how pi approaches security.

\medskip
\noindent\textbf{Status:} case study plus the decisions it forced.

\section*{How it started: building to understand}

This began as a way to understand how pi handles permissions: a tool-call
\textbf{sieve gate}, built to learn the domain rather than to ship a product.

It intercepts each tool call, runs it through an ordered \emph{rejection} sieve
(auto-deny with a teaching reason), then an ordered \emph{approval} sieve
(auto-allow), and sends anything left to a human checkpoint with a plain-language
summary. It hooks pi's \code{tool\_call} event and works: it denies, approves,
and prompts on real commands.

\section*{The pivot: subagents don't share the parent's extension context}

Checking whether the gate would cover \textbf{subagent} activity revealed that it
usually does not. A subagent runs in its own session, and that session only
enforces the extensions \textbf{it} loaded. The parent's gate does not
automatically extend to a child.

That changed the question from ``what rules should the gate apply?'' to:

\begin{pullquote}
Can a security gate in pi even be \emph{guaranteed to run} on the work being done?
\end{pullquote}

Answering it required reading source.

\section*{The investigation}

Read at depth: three subagent extensions (\code{@quintinshaw/pi-dynamic-workflows},
\code{nicobailon/pi-subagents}, \code{tintinweb/pi-subagents}),
\code{@gotgenes/pi-permission-system}'s subagent integration, and the pi harness
itself (the tool-call dispatch path and the resource loader that decides which
extensions a session loads). Three findings.

\subsection*{1. pi's tool boundary is a cooperative, mutable chain}

pi exposes one extension-facing seam at the tool boundary
(\code{agent-session.js:214}): a \code{tool\_call} hook before execution, a
\code{tool\_result} hook after. Dispatch (\code{runner.emitToolCall}) walks
\textbf{all bound extensions in load order}. A handler may return
\code{\{ block, reason \}} or \textbf{mutate the tool input in place with no
re-validation}, and the \textbf{first} block wins. Two properties follow:

\begin{itemize}
  \item \textbf{Enforcement is per-session and opt-in.} The handler set is
    whatever \emph{that} session bound. A session that did not load the gate has
    no gate.
  \item \textbf{The gate is a peer, not an authority.} It shares a mutable input
    with other extensions, order-dependent, so a sibling loaded \emph{after} it
    can rewrite a command the gate already approved.
\end{itemize}

\subsection*{2. Whether a subagent loads the gate varies wildly by extension}

The three answer ``does the parent's policy run in the child?'' differently:

\begin{itemize}
  \item \textbf{\code{dynamic-workflows}: no.} Children are created with
    \code{noExtensions: true} hard-coded; the gate never loads. Tool granting is
    whole-tool allow/deny at spawn (\code{applyToolPolicy}), no command-level
    nuance. It leans on git-worktree isolation, which is damage confinement, not
    access control.
  \item \textbf{\code{nicobailon}: only inside one ecosystem.} It is
    \emph{subprocess}-based, the only one with a real process boundary, but
    exposes no hooks a general policy can bind to. It targets
    \code{pi-permission-system} and misses env wiring
    (\code{PI\_SUBAGENT\_PARENT\_SESSION}) others would need. Outside that
    ecosystem there is no viable integration path.
  \item \textbf{\code{tintinweb}: yes, by default, and scopably.} It loads host
    extensions into children (\code{extensions: true}) and calls
    \code{bindExtensions}, so the gate runs per-call; loading is declaratively
    scopable (\code{extensions: [stacia-sieve-gate]}). It stays in-process (no
    boundary) and reloads everything per child by default, but it is the only one
    that gives a policy real purchase inside subagents.
\end{itemize}

Same policy, unchanged: fully enforced, partially enforced, or entirely absent,
depending on which subagent extension is in use and how it was configured.

\subsection*{3. \code{noExtensions} is a scalpel, which worsens the inconsistency}

\code{noExtensions: true} does not mean ``no extensions in the child''
(\code{resource-loader.js:267}). It drops only the auto-discovered set;
\code{additionalExtensionPaths} and inline \code{extensionFactories} still load.
So \code{noExtensions: true} plus \code{extensionFactories: [gate]} loads
\emph{only} the gate. The capability to guarantee a child is gated exists in the
harness. It is not wired to any consistent, first-class control, so each subagent
extension reinvents or ignores it.

\section*{The conclusion: a structural flaw}

The finding is about pi, not any one extension.

\textbf{Security in pi is a second-class element, an extension composed with peer
extensions, so it can never be guaranteed structurally.} The precise failure is
not ``pi can't stop every attack'' (no in-process model can) but:

\begin{pullquote}
pi cannot guarantee the security layer is even \textbf{present and authoritative}
on a given session.
\end{pullquote}

The gate can be unloaded (\code{isolated} / \code{noExtensions}; the same
\code{noExtensions} is also the scalpel of finding 3), it competes with peers
(load order), and its decisions can be mutated after approval. A security posture
built on pi is defeatable by \emph{configuration} and by \emph{siblings}.

Which risks are pi's and which are common to the paradigm:

\begin{itemize}
  \item \textbf{pi-specific (the differentiators).} A sibling handler mutating an
    approved input, and a child spawned that never loads the gate. Both come from
    the cooperative, opt-in-per-session model. A harness with a single
    always-present authoritative decision point removes both.
  \item \textbf{Paradigm-inherent (shared with any in-process agent, including
    Claude Code).} A Turing-complete shell as an escape hatch, two-step laundering
    with no information-flow control, human-in-the-loop fatigue and broad-grant
    abuse, and rules stored in agent-editable files. These are present in any
    harness and close only with a real boundary (OS sandbox, separate user,
    credential isolation, information-flow control).
\end{itemize}

The differentiator is entirely pi-specific and architectural. pi does not need an
opinion on \emph{how} to implement security. It lacks a \textbf{first-class,
unified substrate}: a harness-owned, always-present, authoritative permission
decision point that security \emph{and} ordinary extensions bind to as policy,
instead of racing as equal peers over a mutable input. Without it, security on pi
is structurally unguaranteeable, and the burden of closing that falls on
userspace.

This is a potential death-knell for adopting pi. The paradigm-inherent vectors
are not pi's fault. The structural one, no guarantee the gate is present and
final, is, and it cannot be patched around from an extension.

\section*{Decisions and mitigations}

The pragmatic path on pi:

\begin{itemize}
  \item \textbf{Adopt \code{tintinweb/pi-subagents}}, the only substrate where the
    gate runs per-call in subagents, scoped to
    \code{extensions: [stacia-sieve-gate]} to keep blast radius small.
  \item \textbf{Forward subagent fall-throughs to the parent UI} via a
    \code{Symbol.for()} broker (child sessions have no UI), with a
    narrowly-scoped grant cache to blunt fatigue.
  \item \textbf{Gate the spawn}, denying \code{Agent} calls that would run ungated
    (\code{isolated: true}, opt-out \code{subagent\_type}s), since every bypass
    the parent can reach is itself a tool call it sees. This blocks
    \emph{unintentional} ungated spawns --- a naive request or a careless config.
    It does \textbf{not} stop an adversarial sibling: per finding 1, a handler
    loaded after the gate can set \code{isolated: true} back \emph{after}
    approval, and the ungated child spawns anyway. The spawn gate would hold only
    if the gate were guaranteed to run last in the chain (validating the final
    input) --- exactly the load-order authority pi cannot grant.
  \item \textbf{Protect the preconditions}, treating writes to
    \code{.pi/agents/**}, \code{\textasciitilde/.pi/agent/extensions/**}, pi
    settings/auth, and pi-re-invoking \code{bash} as top-sensitivity.
\end{itemize}

Every item is userspace mitigation that raises cost and coverage; none is a
structural guarantee. Uncloseable under this design without a harness change or a
real boundary: post-approval mutation, the shell-breadth escape hatch, and the
absence of information-flow control.

\section*{Evidence index}

\begin{itemize}
  \item pi tool hooks / dispatch: \code{agent-session.js:214}
    (\code{beforeToolCall} to \code{emitToolCall}); \code{runner.emitToolCall}
    (load order, first-block short-circuit, mutable input, no re-validation).
  \item \code{noExtensions} semantics: \code{resource-loader.js:267};
    \code{additionalExtensionPaths} and \code{extensionFactories} survive it.
  \item dynamic-workflows: \code{agent.ts} (\code{noExtensions: true};
    \code{applyToolPolicy}; \code{DEFAULT\_EXCLUDED\_SUBAGENT\_TOOLS});
    \code{session.resourceLoader} reachable only to programmatic embedders.
  \item tintinweb: default \code{extensions: true} (\code{agent-types.ts},
    \code{default-agents.ts}); loader-level \code{extensionsOverride};
    \code{isolated} is an \code{Agent} tool param (\code{index.ts}) forcing
    \code{extensions === false} (\code{agent-runner.ts:524});
    \code{bindExtensions} fires \code{session\_start}; child has no UI
    (\code{ctx.hasUI === false}).
  \item nicobailon: subprocess model; missing \code{PI\_SUBAGENT\_PARENT\_SESSION};
    \code{tools:} CSV allowlist.
  \item pi-permission-system: native \code{@gotgenes/pi-subagents} integration
    (process-global registry plus ask-forwarding); two-layer visibility/policy
    model.
\end{itemize}

\end{document}
